% GENERATED FILE. Edit canonical lessons, then run npm run generate:books.
% canonical-source-hash: fnv1a64:cfcd04d009163042
% canonical-lessons: IT-C219-uguale, IT-C219-simile, IT-C219-falso, IT-C219-gratuito, IT-C219-moderno

\chapter{Equal, Similar, False, Free of charge, Modern}
\label{ch:it-equal-similar-false-free-of-charge-modern}
\hlchaptermodality{219}

\begin{chapteropening}
\textbf{By the end of this chapter:} \emph{I can say equal, similar, false, free of charge and modern.}

Complete the last lesson of chapter 219: I can say equal, similar, false, free of charge and modern.
\end{chapteropening}

\section[uguale]{uguale — equal, the same}

\label{lesson:IT-C219-uguale}



\begin{quote}
Before the new one: say the Italian for calm, then the Italian for silent.
\end{quote}

\subsection*{You'll want to know: uguale}
\textbf{uguale} — \textquotedblleft{}equal, the same\textquotedblright{}.

\begin{cousinweb}
\textbf{uguale} $\leftarrow$ Latin \emph{aequalis}, \textquotedblleft{}equal, level\textquotedblright{} (\emph{aequalem}) — inherited. (EtymDB 2.1)
\end{cousinweb}

\begin{grammarlens}[title={What you've built}]
One more word for this chapter.
\end{grammarlens}

\subsection*{Your turn}
\begin{itemize}
\raggedright
  \item \emph{Say it:} \emph{uguale}
  \item \emph{Say it:} \emph{uguale}, once more
  \item \emph{Say it:} say \emph{uguale}
  \item \emph{Recall:} say the Italian for calm, then the Italian for silent, then say \emph{uguale} again
\end{itemize}

\begin{tcolorbox}[breakable,colback=teal!4,colframe=teal!35!black,title={Before you move on}]
What does uguale mean? (\textquotedblleft{}Equal, the same\textquotedblright{}.) Say it once more.
\end{tcolorbox}
\section[simile]{simile — similar}

\label{lesson:IT-C219-simile}



\begin{quote}
Before the new one: say the Italian for silent, then the Italian for equal.
\end{quote}

\subsection*{You'll want to know: simile}
\textbf{simile} — \textquotedblleft{}similar\textquotedblright{}.

\begin{cousinweb}
\textbf{simile} $\leftarrow$ Latin \emph{similis}, \textquotedblleft{}like\textquotedblright{} — inherited. (EtymDB 2.1)
\end{cousinweb}

\begin{grammarlens}[title={What you've built}]
One more word for this chapter.
\end{grammarlens}

\subsection*{Your turn}
\begin{itemize}
\raggedright
  \item \emph{Say it:} \emph{simile}
  \item \emph{Say it:} \emph{simile}, once more
  \item \emph{Say it:} say \emph{simile}
  \item \emph{Recall:} say the Italian for silent, then the Italian for equal, then say \emph{simile} again
\end{itemize}

\begin{tcolorbox}[breakable,colback=teal!4,colframe=teal!35!black,title={Before you move on}]
What does simile mean? (\textquotedblleft{}Similar\textquotedblright{}.) Say it once more.
\end{tcolorbox}
\section[falso]{falso — false}

\label{lesson:IT-C219-falso}



\begin{quote}
Before the new one: say the Italian for equal, then the Italian for similar.
\end{quote}

\subsection*{You'll want to know: falso}
\textbf{falso} — \textquotedblleft{}false\textquotedblright{}.

\begin{cousinweb}
\textbf{falso} $\leftarrow$ Latin \emph{falsus}, \textquotedblleft{}deceived, false\textquotedblright{}, from \emph{fallere}, \textquotedblleft{}to deceive\textquotedblright{} — inherited. (EtymDB 2.1)
\end{cousinweb}

\begin{grammarlens}[title={What you've built}]
One more word for this chapter.
\end{grammarlens}

\subsection*{Your turn}
\begin{itemize}
\raggedright
  \item \emph{Say it:} \emph{falso}
  \item \emph{Say it:} \emph{falso}, once more
  \item \emph{Say it:} say \emph{falso}
  \item \emph{Recall:} say the Italian for equal, then the Italian for similar, then say \emph{falso} again
\end{itemize}

\begin{tcolorbox}[breakable,colback=teal!4,colframe=teal!35!black,title={Before you move on}]
What does falso mean? (\textquotedblleft{}False\textquotedblright{}.) Say it once more.
\end{tcolorbox}
\section[gratuito]{gratuito — free of charge}

\label{lesson:IT-C219-gratuito}



\begin{quote}
Before the new one: say the Italian for similar, then the Italian for false.
\end{quote}

\subsection*{You'll want to know: gratuito}
\textbf{gratuito} — \textquotedblleft{}free of charge\textquotedblright{}.

\begin{cousinweb}
\textbf{gratuito} $\leftarrow$ Latin \emph{gratuitus}, \textquotedblleft{}free, done for nothing\textquotedblright{}, from \emph{gratia}, \textquotedblleft{}favour\textquotedblright{} — a learned borrowing. (EtymDB 2.1)
\end{cousinweb}

\begin{grammarlens}[title={What you've built}]
One more word for this chapter.
\end{grammarlens}

\subsection*{Your turn}
\begin{itemize}
\raggedright
  \item \emph{Say it:} \emph{gratuito}
  \item \emph{Say it:} \emph{gratuito}, once more
  \item \emph{Say it:} say \emph{gratuito}
  \item \emph{Recall:} say the Italian for similar, then the Italian for false, then say \emph{gratuito} again
\end{itemize}

\begin{tcolorbox}[breakable,colback=teal!4,colframe=teal!35!black,title={Before you move on}]
What does gratuito mean? (\textquotedblleft{}Free of charge\textquotedblright{}.) Say it once more.
\end{tcolorbox}
\section[moderno]{moderno — modern}

\label{lesson:IT-C219-moderno}



\begin{quote}
Before the new one: say the Italian for false, then the Italian for free of charge.
\end{quote}

\subsection*{You'll want to know: moderno}
\textbf{moderno} — \textquotedblleft{}modern\textquotedblright{}.

\begin{cousinweb}
\textbf{moderno} $\leftarrow$ Late Latin \emph{modernus}, \textquotedblleft{}of just now\textquotedblright{}, from \emph{modo}, \textquotedblleft{}just now\textquotedblright{} — a learned borrowing. EtymDB 2.1 has no row for \emph{moderno}; it derives French \emph{moderne} and Catalan \emph{modern} from \emph{modernus}.
\end{cousinweb}

\begin{grammarlens}[title={What you've built}]
One more word for this chapter.
\end{grammarlens}

\subsection*{Your turn}
\begin{itemize}
\raggedright
  \item \emph{Say it:} \emph{moderno}
  \item \emph{Say it:} \emph{moderno}, once more
  \item \emph{Say it:} say \emph{moderno}
  \item \emph{Recall:} say the Italian for false, then the Italian for free of charge, then say \emph{moderno} again
\end{itemize}

\begin{tcolorbox}[breakable,colback=teal!4,colframe=teal!35!black,title={Before you move on}]
What does moderno mean? (\textquotedblleft{}Modern\textquotedblright{}.) Say it once more.
\end{tcolorbox}
